%Version 3.1 December 2024
% See section 11 of the User Manual for version history
%
%%%%%%%%%%%%%%%%%%%%%%%%%%%%%%%%%%%%%%%%%%%%%%%%%%%%%%%%%%%%%%%%%%%%%%
%%                                                                 %%
%% Please do not use \input{...} to include other tex files.       %%
%% Submit your LaTeX manuscript as one .tex document.              %%
%%                                                                 %%
%% All additional figures and files should be attached             %%
%% separately and not embedded in the \TeX\ document itself.       %%
%%                                                                 %%
%%%%%%%%%%%%%%%%%%%%%%%%%%%%%%%%%%%%%%%%%%%%%%%%%%%%%%%%%%%%%%%%%%%%%

%%\documentclass[referee,sn-basic]{sn-jnl}% referee option is meant for double line spacing

%%=======================================================%%
%% to print line numbers in the margin use lineno option %%
%%=======================================================%%

%%\documentclass[lineno,pdflatex,sn-basic]{sn-jnl}% Basic Springer Nature Reference Style/Chemistry Reference Style

%%=========================================================================================%%
%% the documentclass is set to pdflatex as default. You can delete it if not appropriate.  %%
%%=========================================================================================%%

%%\documentclass[sn-basic]{sn-jnl}% Basic Springer Nature Reference Style/Chemistry Reference Style

%%Note: the following reference styles support Namedate and Numbered referencing. By default the style follows the most common style. To switch between the options you can add or remove “Numbered” in the optional parenthesis. 
%%The option is available for: sn-basic.bst, sn-chicago.bst%  
 
%%\documentclass[pdflatex,sn-nature]{sn-jnl}% Style for submissions to Nature Portfolio journals
%%\documentclass[pdflatex,sn-basic]{sn-jnl}% Basic Springer Nature Reference Style/Chemistry Reference Style
\documentclass[pdflatex,sn-mathphys-num]{sn-jnl}% Math and Physical Sciences Numbered Reference Style
%%\documentclass[pdflatex,sn-mathphys-ay]{sn-jnl}% Math and Physical Sciences Author Year Reference Style
%%\documentclass[pdflatex,sn-aps]{sn-jnl}% American Physical Society (APS) Reference Style
%%\documentclass[pdflatex,sn-vancouver-num]{sn-jnl}% Vancouver Numbered Reference Style
%%\documentclass[pdflatex,sn-vancouver-ay]{sn-jnl}% Vancouver Author Year Reference Style
%%\documentclass[pdflatex,sn-apa]{sn-jnl}% APA Reference Style
%%\documentclass[pdflatex,sn-chicago]{sn-jnl}% Chicago-based Humanities Reference Style

%%%% Standard Packages
%%<additional latex packages if required can be included here>

\usepackage{graphicx}%
\usepackage{multirow}%
\usepackage{amsmath,amssymb,amsfonts}%
\usepackage{amsthm}%
\usepackage{mathrsfs}%
\usepackage[title]{appendix}%
\usepackage{xcolor}%
\usepackage{textcomp}%
\usepackage{manyfoot}%
\usepackage{booktabs}%
\usepackage{algorithm}%
\usepackage{algorithmicx}%
\usepackage{algpseudocode}%
\usepackage{listings}%
\usepackage{bm}%
%%%%

%%%%%=============================================================================%%%%
%%%%  Remarks: This template is provided to aid authors with the preparation
%%%%  of original research articles intended for submission to journals published 
%%%%  by Springer Nature. The guidance has been prepared in partnership with 
%%%%  production teams to conform to Springer Nature technical requirements. 
%%%%  Editorial and presentation requirements differ among journal portfolios and 
%%%%  research disciplines. You may find sections in this template are irrelevant 
%%%%  to your work and are empowered to omit any such section if allowed by the 
%%%%  journal you intend to submit to. The submission guidelines and policies 
%%%%  of the journal take precedence. A detailed User Manual is available in the 
%%%%  template package for technical guidance.
%%%%%=============================================================================%%%%

%% as per the requirement new theorem styles can be included as shown below
\theoremstyle{thmstyleone}%
\newtheorem{theorem}{Theorem}%  meant for continuous numbers
%%\newtheorem{theorem}{Theorem}[section]% meant for sectionwise numbers
%% optional argument [theorem] produces theorem numbering sequence instead of independent numbers for Proposition
\newtheorem{proposition}[theorem]{Proposition}% 
%%\newtheorem{proposition}{Proposition}% to get separate numbers for theorem and proposition etc.

\theoremstyle{thmstyletwo}%
\newtheorem{example}{Example}%
\newtheorem{remark}{Remark}%

\theoremstyle{thmstylethree}%
\newtheorem{definition}{Definition}%

\raggedbottom
%%\unnumbered% uncomment this for unnumbered level heads

\begin{document}

\title[Article Title]{Differentiable Physics for Physics-Informed Kirigami Inverse Form Finding}

%%=============================================================%%
%% GivenName	-> \fnm{Joergen W.}
%% Particle	-> \spfx{van der} -> surname prefix
%% FamilyName	-> \sur{Ploeg}
%% Suffix	-> \sfx{IV}
%% \author*[1,2]{\fnm{Joergen W.} \spfx{van der} \sur{Ploeg} 
%%  \sfx{IV}}\email{iauthor@gmail.com}
%%=============================================================%%

\author[1,2]{\fnm{Julien} \sur{Kloers}}\email{jk3804@princeton.edu}

\author[2]{\fnm{Sigrid} \sur{Adriaenssens}}

\author[3]{\fnm{Rafael} \sur{Pastrana}}

\affil[1]{\orgname{École des Ponts}, \orgaddress{\street{8 av. Blaise Pascal}, \city{Champs-sur-Marne}, \postcode{77420}, \country{France}}}

\affil[2]{\orgdiv{Form-Finding Lab}, \orgname{Princeton University}, \orgaddress{\street{59 Olden Street}, \city{Princeton}, \postcode{08544}, \state{NJ}, \country{USA}}}

\affil[3]{\orgdiv{Princeton AI2 Lab}, \orgname{Princeton University}, \orgaddress{\street{41 William Street}, \city{Princeton}, \postcode{08540}, \state{NJ}, \country{USA}}}

%%==================================%%
%% Sample for unstructured abstract %%
%%==================================%%

\abstract{The abstract serves both as a general introduction to the topic and as a brief, non-technical summary of the main results and their implications. Authors are advised to check the author instructions for the journal they are submitting to for word limits and if structural elements like subheadings, citations, or equations are permitted.}

%%================================%%
%% Sample for structured abstract %%
%%================================%%

% \abstract{\textbf{Purpose:} The abstract serves both as a general introduction to the topic and as a brief, non-technical summary of the main results and their implications. The abstract must not include subheadings (unless expressly permitted in the journal's Instructions to Authors), equations or citations. As a guide the abstract should not exceed 200 words. Most journals do not set a hard limit however authors are advised to check the author instructions for the journal they are submitting to.
% 
% \textbf{Methods:} The abstract serves both as a general introduction to the topic and as a brief, non-technical summary of the main results and their implications. The abstract must not include subheadings (unless expressly permitted in the journal's Instructions to Authors), equations or citations. As a guide the abstract should not exceed 200 words. Most journals do not set a hard limit however authors are advised to check the author instructions for the journal they are submitting to.
% 
% \textbf{Results:} The abstract serves both as a general introduction to the topic and as a brief, non-technical summary of the main results and their implications. The abstract must not include subheadings (unless expressly permitted in the journal's Instructions to Authors), equations or citations. As a guide the abstract should not exceed 200 words. Most journals do not set a hard limit however authors are advised to check the author instructions for the journal they are submitting to.
% 
% \textbf{Conclusion:} The abstract serves both as a general introduction to the topic and as a brief, non-technical summary of the main results and their implications. The abstract must not include subheadings (unless expressly permitted in the journal's Instructions to Authors), equations or citations. As a guide the abstract should not exceed 200 words. Most journals do not set a hard limit however authors are advised to check the author instructions for the journal they are submitting to.}

\keywords{keyword1, Keyword2, Keyword3, Keyword4}

%%\pacs[JEL Classification]{D8, H51}

%%\pacs[MSC Classification]{35A01, 65L10, 65L12, 65L20, 65L70}

\maketitle

\section{Introduction}\label{sec1}

Kirigami

Physique différentiable, JAX

Hypotheses : 
Rigid planar quadrangular deployable kirigami. 
Les faces sont rigides et coplanaiaires pendant tout le développement. 
On étudie en particulier le 


Plan de l'article

\section{Methods}\label{secM}

\subsection{The kirigami tessellation}

We model a planar kirigami metamaterial as an assembly of $n$ rigid polygonal tiles cut from a flat sheet and connected to one another, at their corners, by slender elastic hinges. Deformation is assumed to localise entirely in the ligaments: the panels themselves are treated as rigid. This
reduction is appropriate in the regime where the panels are much stiffer than the thin connecting ligaments, and it collapses a high-dimensional continuum problem onto at most $3n$ degrees of freedom. A tessellation is called deployable, if it has exactly one degree of freedom. 

\paragraph{Topology.}
The connectivity is fixed by the cut pattern and does not change during deformation. A cut pattern is the arrangement of slits that decomposes a continuous sheet into rigid panels joined by intact ligaments, thereby fixing both the connectivity and the flat reference geometry of the kirigami. We generate it from the closed-state RDPQK construction of Dang~\emph{et al.}, laid out on an $M \times N$ grid of unit cells. From this pattern, we derive once and for all the purely combinatorial data of the assembly: the set of faces
$\{P_\alpha\}_{\alpha=1}^{n}$, the corner (node) incidences of each face, and the set of \emph{hinges} $\mathcal{H}$, where each hinge $h=(i,k)\in\mathcal{H}$
identifies a pair of face corners that an uncut strip of material connects. 

\paragraph{Geometry and design parameters.}
Each rigid face $P_\alpha$ is described in a centroidal representation by its
centroid $\bm{c}_\alpha \in \mathbb{R}^2$ together with a set of
centroid-to-corner vectors $\{\bm{d}_{\alpha j}\}_j$; the corner positions in the
reference (flat, undeployed) configuration are
$\bm{x}_{\alpha j}^{0} = \bm{c}_\alpha + \bm{d}_{\alpha j}$. This reference state is
produced from the trainable design parameters
\begin{equation}
  \bm{\theta} = \bigl(\, \bm{z},\ \bm{\ell},\ \bm{w}_{\mathrm{lig}}\,\bigr),
  \qquad
  \bm{r} = \sigma(\bm{z}) \in (0,1)^{\lvert\mathrm{cuts}\rvert},
  \label{eq:design-params}
\end{equation}
where the per-cut logits $\bm{z}$ map through the sigmoid $\sigma$ to \emph{aspect
ratios} $\bm{r}$, each $r$ fixing the relative placement of a cut's two endpoints;
the ordered logits $\bm{\ell}$ map through a cumulative-softmax to the tangential
positions $\bm{s}$ of the boundary cuts (\emph{boundary sliders}); and
$\bm{w}_{\mathrm{lig}}$ collects the per-hinge ligament widths. Given the induced
$(\bm{r},\bm{s})$, the interior cut vertices are obtained as the solution of a
linear collinearity system
\begin{equation}
  A(\bm{r})\,\bm{v} = \bm{b}(\bm{s}),
\end{equation}
which places every interior vertex on the line through its neighbours with the
prescribed aspect ratio, while boundary vertices are pinned to the slider
positions. The map $\bm{\theta} \mapsto \{\bm{c}_\alpha, \bm{d}_{\alpha j}\}$ is
smooth and the resulting flat tessellation is guaranteed to be
non-self-intersecting: the cut aspect ratios yield parallelogram voids, and the
boundary sliders are kept strictly ordered (the cumulative-softmax
parameterisation) so that the outline stays convex, satisfying the hypotheses of
Tutte's embedding theorem for any admissible $\bm{\theta}$. Validity of the flat
state is therefore \emph{built in}, not enforced a posteriori.

Alongside the shape parameters, each hinge $h$ carries two geometric attributes
that are constants of the design and enter the mechanics as material descriptors:
its rest orientation, summarised by a cut angle $\beta_h$ (a function of the
neighbouring $\bm{r}$), and its ligament width $w_{\mathrm{lig},h}$, which sets the
local compliance. These are fixed by the flat construction and do not evolve during
deployment.

\subsection{Kinematics}

Because the panels are rigid, the configuration of the whole structure is
specified by three degrees of freedom per face,
\begin{equation}
  \bm{q}_\alpha = (\bm{u}_\alpha,\ \phi_\alpha) \in \mathbb{R}^2 \times \mathbb{S}^1,
  \qquad
  \bm{q} = (\bm{q}_1,\dots,\bm{q}_{n}) \in \mathbb{R}^{3n},
\end{equation}
a rigid-body translation $\bm{u}_\alpha$ and rotation $\phi_\alpha$ of face
$P_\alpha$ about its centroid. The displacement of corner $j$ of face $\alpha$
follows exactly from the rigid-body map,
\begin{equation}
  \Delta\bm{x}_{\alpha j}(\bm{q}_\alpha)
    = \bm{u}_\alpha + \bigl(R(\phi_\alpha) - I\bigr)\,\bm{d}_{\alpha j},
\end{equation}
where $R(\phi)$ is the planar rotation matrix; the rotation is kept fully
nonlinear (no small-angle linearisation of the panel motion), so that the model
remains valid through the large panel rotations that characterise deployment.

\subsection{Hinge deformation measures}

The mechanical state of a hinge is captured by the relative motion of the two
faces it connects, expressed in a frame that is invariant to rigid-body motion of
the assembly. For hinge $h=(i,k)$ --- joining faces $i$ and $k$ --- with reference
bond vector $\bm{b}_h^{0} = \bm{x}_k^{0} - \bm{x}_i^{0}$, we remove the mean
rotation of the two adjacent faces (a corotational, total-Lagrangian frame) and
read off three scalar deformation measures,
\begin{equation}
  \bm{\varepsilon}_h(\bm{q}) = \bigl(a_h,\ s_h,\ \vartheta_h\bigr),
  \label{eq:hinge-strain}
\end{equation}
namely an axial (stretch) measure $a_h$, a shear measure $s_h$, and the relative
rotation $\vartheta_h = \phi_k - \phi_i$. By construction $\bm{\varepsilon}_h$
vanishes under any rigid-body motion of the structure, including arbitrarily large
global rotations, and depends only on the \emph{current} configuration relative to
the reference geometry --- the ligament is treated as memoryless. 

\begin{figure}
    \centering
    \includegraphics[width=1\linewidth]{Images/pipeline_overview.png}
    \caption{Enter Caption}
    \label{fig:placeholder}
\end{figure}

\subsection{Energy functional}

We equip each hinge with a stored-energy function $W$ of its own deformation
measures and design attributes,
\begin{equation}
  W_h = W\!\bigl(\bm{\varepsilon}_h;\ \bm{g}_h\bigr),
  \qquad
  \bm{g}_h = \bigl(w_{\mathrm{lig},h},\ \beta_h\bigr),
  \label{eq:hinge-energy}
\end{equation}
where $\bm{g}_h$ collects the per-hinge design attributes (ligament width and cut
angle), and take the total elastic energy of the assembly to be the sum of
independent hinge contributions,
\begin{equation}
  U(\bm{q};\,\bm{\theta}) = \sum_{h\in\mathcal{H}}
    W\!\bigl(\bm{\varepsilon}_h(\bm{q});\ \bm{g}_h\bigr).
  \label{eq:total-energy}
\end{equation}
We deliberately leave the constitutive form of $W$ abstract: any objective,
frame-invariant function of $\bm{\varepsilon}=(a,s,\vartheta)$, parameterised by
the hinge descriptors $\bm{g}_h$, is admissible, and different choices --- from a
simple quadratic spring to a data-driven surrogate calibrated against
high-fidelity simulation --- can be substituted without altering the surrounding
kinematics or topology. The only requirements are that $W$ be non-negative,
minimised at the undeformed state, and differentiable in both its kinematic
arguments and its design attributes.

\subsection{Equilibrium}

Deployment and loading are imposed through boundary conditions on the face
degrees of freedom: a subset of face DOFs is prescribed (Dirichlet, the clamped
panels) and generalised forces $\bm{f}$ are applied to others (Neumann). The
mechanical state at load factor $\lambda\in[0,1]$ is the minimiser of the total
potential energy over the free degrees of freedom,
\begin{equation}
  \bm{q}^\star(\lambda;\,\bm{\theta})
    = \arg\min_{\bm{q}\,\in\,\mathcal{C}}
      \ \Pi(\bm{q};\,\bm{\theta},\lambda),
  \qquad
  \Pi = U(\bm{q};\,\bm{\theta}) - \lambda\,\bm{f}^\top \bm{q},
  \label{eq:equilibrium}
\end{equation}
where $\mathcal{C}$ is the affine space of configurations satisfying the
prescribed DOFs. The load is applied incrementally in $\lambda$, each increment
warm-started from the previous equilibrium, so that the model traces a
quasi-static deformation path from the flat closed sheet to its deployed
configuration. Everything downstream of \eqref{eq:equilibrium} --- the equilibrium
map $\bm{\theta} \mapsto \bm{q}^\star$ --- is differentiable with respect to the
design parameters $\bm{\theta}$, which is what makes the object amenable to
inverse design.


\subsection{Differentiable physics pipeline}

In the closed inverse-design setting the trainable parameters are those of
\eqref{eq:design-params}, $\bm{\theta}=(\bm{z},\bm{\ell},\bm{w}_{\mathrm{lig}})$:
per-cut logits $\bm{z}$ mapped to Tutte aspect ratios
$\bm{r}=\sigma(\bm{z})\in(0,1)$, ordered boundary-slider logits $\bm{\ell}$, and
per-hinge ligament widths $\bm{w}_{\mathrm{lig}}$. A differentiable forward map
$\mathcal{G}_{\bm{\theta}}$ (the flat RDPQK closed-sheet builder) assembles
$\bm{r}$ and the boundary sliders into a sparse collinearity linear system whose
solution gives the flat panel vertices; non-self-intersection holds by
construction. Deployment is the static solve
\begin{equation}
  \bm{q}^\star(\bm{\theta})
    = \arg\min_{\bm{q}\,\in\,\mathcal{C}}\,\Pi(\bm{q};\,\bm{\theta},\lambda),
  \qquad
  \Pi = \sum_{h\in\mathcal{H}} W\!\bigl(\bm{\varepsilon}_h(\bm{q});\,\bm{g}_h\bigr)
        - W_{\mathrm{ext}},
  \label{eq:pipeline-solve}
\end{equation}
with external work $W_{\mathrm{ext}}=\lambda\,\bm{f}^\top\bm{q}$ and the condensed
hinge energy $W$ the learned (frozen) $\tanh$-MLP surrogate of
\eqref{eq:surrogate} below; its automatic-differentiation gradient
$\nabla_{\bm{\varepsilon}} W=(F_a,F_s,M_\vartheta)$ supplies the internal
generalised forces and $\nabla^2_{\bm{\varepsilon}} W$ the tangent stiffness.

The minimisation uses jaxopt L-BFGS under incremental load stepping (in $\lambda$)
unrolled with \texttt{jax.lax.scan}. Gradients propagate through the argmin by the
implicit function theorem rather than through the solver iterates,
\begin{equation}
  \frac{\partial \bm{q}^\star}{\partial \bm{\theta}}
    = -\,H^{-1}\,\frac{\partial^2 \Pi}{\partial \bm{q}\,\partial \bm{\theta}},
  \qquad
  H = \nabla^2_{\bm{q}}\,\Pi,
\end{equation}
exposed as a \texttt{custom\_vjp}. Since plastic softening in the surrogate can
render $H$ indefinite (a vanishing design gradient), the backward solve is
Tikhonov-regularised, $H\!\leftarrow\!H+\tau I$
(\texttt{physics.backward\_reg}), restoring a descent direction. Finally the
end-to-end objective --- a Chamfer distance between the deployed centroids
$\{\bm{c}_\alpha(\bm{q}^\star)\}$ and the target shape $\mathcal{T}$, plus a
physical-stability penalty --- is minimised with Adam through a single
\texttt{jax.value\_and\_grad}, so gradients flow
$\mathcal{L}\!\to\!\bm{q}^\star\!\to\!\Pi\!\to\!\bm{\theta}$ back to the cut
aspect ratios, boundary sliders, and ligament widths jointly.

\subsection{Condensed hinge-energy surrogate}

% Buckling ?

Each kirigami hinge (a thin ligament joining two rigid tiles) is replaced by a
single learned \emph{stored-energy function} $W(\bm{\varepsilon};\bm{g})$
[N$\cdot$mm]. It condenses a full 3D elasto-plastic simulation into a scalar that
the deployment solver can minimise, and whose gradient is the internal generalised
force the solver balances. The relative motion of the two connected faces is the
minimal rigid-body-invariant coordinate $\bm{\varepsilon}=(a,\,s,\,\vartheta)$ of
\eqref{eq:hinge-strain}: the corotated axial, shear, and relative-rotation of the
inter-tile bond. The design fixes the per-hinge attributes
$\bm{g}=(w_{\mathrm{lig}},\,\beta\,[,\rho_f])$ --- ligament width, cut angle, and
the cut-tip fillet ratio --- with both $\beta$ and the bond direction
$\hat{\bm{e}}_\parallel$ deterministic functions of the void aspect ratios
$\bm{r}$ and boundary sliders. The geometry $\bm{g}(\bm{\theta})$ is differentiable
and threaded as an \emph{explicit} solver input, so the shape loss $\mathcal{L}$
(target Chamfer $+$ stability penalty on $D$) is minimised through the physics,
including the design's effect on $\beta(\bm{\theta})$. With
$\eta:\mathbb{R}^{6}\!\to\mathbb{R}^m$ a $\tanh$ MLP acting on the standardised
features
$\bm{\xi}(\bm{\varepsilon},\bm{g})=[a,s,\vartheta,\log w_{\mathrm{lig}},\beta,\rho_f]$,
\begin{equation}
  \boxed{\,W(\bm{\varepsilon};\bm{g})
    = W_{s}\,\bigl\|\,\eta\bigl(\bm{\xi}(\bm{\varepsilon},\bm{g})\bigr)
      -\eta\bigl(\bm{\xi}(\bm{0},\bm{g})\bigr)\,\bigr\|_2^2\,}.
  \label{eq:surrogate}
\end{equation}
The squared form gives, \emph{by construction} and with no penalty terms, no
spurious negative-energy state, free rigid-body motion, and a smooth tangent
stiffness for the implicit solve. The internal generalised force is the gradient
\begin{equation}
  \nabla_{\bm{\varepsilon}}W=(F_a,\,F_s,\,M_\vartheta),
\end{equation}
and a separate head predicts a continuous ductile-damage margin
$D(\bm{\varepsilon};\bm{g})=\mathrm{softplus}(\cdot)\ge 0$ ($D\!\ge\!1$ =
fracture), kept distinct so the hard failure boundary does not distort the smooth
$W$.

The scale $W_{s}$ and the maps $\eta,\,D$ are fit on a CalculiX RVE dataset by a
variance-normalised Sobolev loss (energy priority via the weight $\gamma$),
\begin{equation}
  \mathcal{L}_{\mathrm{data}}
  =\gamma\,\frac{\|W-W^\ast\|^2}{\sigma_W^2}
  +(1-\gamma)\,\frac{\|\nabla_{\bm{\varepsilon}}W-\bm{F}^\ast\|^2}{\sigma_F^2}
  +w_f\,\|D-D^\ast\|^2,
\end{equation}
so both the energy landscape and the equilibrium-setting forces are matched, split
by geometry to test extrapolation to unseen hinges.

\subsubsection{Finite element modeling of the hinges}

Each ligament hinge is characterized by a three-dimensional finite-element model of a Saint--Venant representative volume element built around the cut, solved with CalculiX. The domain is discretized with quadratic 15-node pentahedral elements (C3D15) and the S235 structural steel is described by a
finite-strain $J_2$ (von~Mises) elastoplastic law with isotropic hardening ($E=210\,\mathrm{GPa}$, $\sigma_y\!\approx\!235\,\mathrm{MPa}$, tangent modulus
$E_t\!\approx\!2.1\,\mathrm{GPa}$), so that yielding, mild hardening, and the
geometrically nonlinear out-of-plane buckling of the ligament are all captured.
The RVE is loaded in a co-rotational frame by the three relative face-to-face
deformation measures $\bm{\varepsilon}=(a,s,\vartheta)$ (axial and shear openings
and the hinge rotation), imposed as kinematic boundary conditions on the two
bounding faces while the far field is Saint--Venant--relaxed. For each geometry
$\bm{g}=(w_{\mathrm{lig}},\beta)$ (ligament width and cut angle) and load state,
the simulation returns the condensed stored energy
$W(\bm{\varepsilon};\bm{g})=\!\int_\Omega \psi\,\mathrm{d}V$, its work-conjugate
generalized forces $\nabla_{\bm{\varepsilon}}W=(F_a,F_s,M_\vartheta)$, the peak
von~Mises stress, and the 99th-percentile equivalent plastic strain
$\varepsilon_p^{p99}$, whose ratio to the ductile failure strain $\varepsilon_f$
defines the damage margin
$D=\varepsilon_p^{p99}/\varepsilon_f$ ($D\!\ge\!1$ marking fracture).

\begin{figure}
    \centering
    \includegraphics[width=1\linewidth]{Images/hinge_layout.png}
    \caption{Enter Caption}
    \label{fig:hinge_layout}
\end{figure}


\begin{figure}
    \centering
    \includegraphics[width=1\linewidth]{Images/energy_landscape.png}
    \caption{Surrogate learned rotation energy landscape of steel hinges}
    \label{fig:placeholder}
\end{figure}

\subsection{Calibration of the surrogate to a physical material}
\label{sec:calibration}

The surrogate of \eqref{eq:surrogate} is only as physical as the constitutive law
handed to the representative volume element that generates its training data. To
build and cut real specimens we therefore retarget the oracle from the reference
S235 steel to the sheet actually used in the experiments --- a $0.5\,\mathrm{mm}$
biaxially-oriented polyethylene terephthalate (PET) film --- through the
calibration chain
\begin{equation}
  \text{coupon}\ \longrightarrow\
  \text{bulk law}\ \longrightarrow\
  \text{RVE (CalculiX)}\ \longrightarrow\
  W(\bm{\varepsilon};\bm{g})\ \longrightarrow\
  \text{deployment},
  \label{eq:calibration-chain}
\end{equation}
in which each arrow is identified or verified separately. The guiding principle is
a strict separation of \emph{material} from \emph{structure}: the bulk law is
identified from plain tensile coupons alone, and the hinge experiments are then
used to \emph{validate} the resulting predictions rather than to fit them. The
only genuinely structural quantity left free is the geometric imperfection
amplitude that seeds the out-of-plane buckling of the ligament, which is not a
material property and cannot be measured on a coupon. This keeps the surrogate a
prediction of the hinge response rather than a curve fit to it.

\paragraph{Material identification.}
The sheet was identified from bulk indicators alone: a measured density
$\rho\simeq1390\,\mathrm{kg\,m^{-3}}$ (excluding PETG, polycarbonate and PMMA), a
broad-band iridescence indicative of molecular orientation, a $180^\circ$ fold
without rupture, and a burn test in which the specimen melted, dripped and drew
into long filaments without sustained burning --- the signature of a thermoplastic
polyester, which also excludes PVC. Under load the film cold-draws through a
propagating neck and becomes birefringent, confirming strain-induced orientation.

\paragraph{Bulk law from tensile coupons.}
Ten rectangular strips (ASTM~D882 geometry, six along the machine direction and
four across it) were pulled to the cold-draw plateau on a $5\,\mathrm{kN}$
single-column frame at $1\,\mathrm{mm\,min^{-1}}$, with tape-tabbed flat grips.
Nine drew stably. The yield stress is direction-independent within scatter
($48.4\,\mathrm{MPa}$ machine direction versus $49.5\,\mathrm{MPa}$ cross
direction, pooled $48.8\pm1.6\,\mathrm{MPa}$), so the film is treated as in-plane
isotropic and modelled with a two-constant \texttt{*ELASTIC} card rather than the
nine-constant orthotropic route. Because the neck localises the deformation, the
engineering plateau is not the material law; the post-yield point is instead
recovered from the measured natural draw ratio $\lambda=A_0/A_{\mathrm{drawn}}$,
\begin{equation}
  \sigma_{\mathrm{true}} = \lambda\,\sigma_{\mathrm{plateau}},
  \qquad
  \varepsilon_{\mathrm{true}} = \ln\lambda,
  \label{eq:draw-anchor}
\end{equation}
which for $\lambda=3.28$ and an engineering plateau of $33\,\mathrm{MPa}$ places a
large-strain anchor at $109\,\mathrm{MPa}$ at a true plastic strain of $1.15$.
The propagating neck also halves the section at constant volume, confirming
plastic incompressibility and hence the admissibility of a $J_2$ flow rule. The
resulting two-point monotonic hardening table (Table~\ref{tab:pet-params})
replaces the bilinear steel law by a curve anchored in the same large-strain
regime the ligament actually experiences. No softening branch is retained: a
ligament loaded in bending forms a crease, not a propagating neck, so the
relevant plastic response is monotonic.

\begin{table}[h!]
\centering
\begin{tabular}{llll}
\toprule
\textbf{Quantity} & \textbf{Symbol} & \textbf{Value} & \textbf{Source} \\
\midrule
Young's modulus        & $E$                     & $3.0\,\mathrm{GPa}$              & literature (see text) \\
Poisson's ratio        & $\nu$                   & $0.40$                           & literature \\
Yield stress           & $\sigma_y$              & $48.8\pm1.6\,\mathrm{MPa}$       & 9 coupons, MD $+$ CD \\
Cold-draw anchor       & $(\sigma,\varepsilon_p)$& $(109\,\mathrm{MPa},\,1.15)$     & Eq.~\eqref{eq:draw-anchor}, $\lambda=3.28$ \\
Fracture strain        & $\varepsilon_{f0}$      & $\geq1.2$                        & lower bound (see text) \\
Sheet thickness        & $t$                     & $0.50\,\mathrm{mm}$              & micrometer \\
Density                & $\rho$                  & $1390\,\mathrm{kg\,m^{-3}}$      & mass\,/\,area\,/\,$t$ \\
\bottomrule
\end{tabular}
\caption{Calibrated PET parameters entering the CalculiX RVE. Elastic constants
are literature values; the plastic table and the failure strain are measured.}
\label{tab:pet-params}
\end{table}

\paragraph{What the coupons cannot deliver.}
Two quantities in Table~\ref{tab:pet-params} are deliberately not taken from the
crosshead data. First, with no extensometer the crosshead displacement contains
machine and grip compliance, which inflates the strain and depresses the apparent
modulus to $\approx1.2\,\mathrm{GPa}$ --- below the entire physical range for PET
--- so a literature value of $3.0\,\mathrm{GPa}$ is used instead. This is
inconsequential for the quantities the surrogate is calibrated against: repeating
a hinge opening simulation at $E=1.2$ and $3.0\,\mathrm{GPa}$ changes the peak
force by $2\%$ while the elastic stiffness scales with $E$, confirming that the
peak is set by the yield stress, the hardening and the failure strain, and that
adopting the compliance-corrupted modulus would merely bake machine compliance
into the material. Second, the specimens were stopped on the draw plateau rather
than broken, so the fracture strain is reported as a lower bound; it is refined
from the hinge tests, where fracture is the quantity of interest. Failure is
modelled as triaxiality-dependent ductile damage,
$D=\varepsilon_p^{p99}/\varepsilon_f(\eta)$ with
$\varepsilon_f(\eta)=\varepsilon_{f0}\exp\!\bigl[-k(\eta-\tfrac13)\bigr]$
--- a Johnson--Cook-like locus normalised so that $\varepsilon_f=\varepsilon_{f0}$
in uniaxial tension --- consistent with the conjugate stress-whitening
bands observed at $\pm45^\circ$ to the load axis: shear-driven yielding on the
planes of maximum shear, not $90^\circ$ crazing, which supports the von~Mises
framework and indicates weak pressure sensitivity. Whitening onset doubles as a
visible, photographable proxy for $D$ and thereby anchors the level of damage that
constitutes functional failure.

\paragraph{Sweeping the deformation space with a one-degree-of-freedom machine.}
A uniaxial testing frame cannot directly impose the three hinge measures
$\bm{\varepsilon}=(a,s,\vartheta)$. The in-plane pair is nonetheless reachable
without any fixture by rotating the specimen in the grips: with $\varphi$ the
angle between the pull direction and the cut, a single tensile pull realises the
ray
\begin{equation}
  (\eta_a,\eta_s)=(\sin\varphi,\ \cos\varphi),
\end{equation}
so that $\varphi=0$ is pure shear along the ligament, $\varphi=90^\circ$ is pure
opening, and intermediate angles are mixed states that additionally sample
different notch triaxialities --- which is precisely what is needed to calibrate
the pair $(\varepsilon_{f0},k)$ of the failure locus. The relative rotation
$\vartheta$, the dominant degree of freedom during deployment, is not accessible
in force control; it is instead validated kinematically, by folding a hinge to a
prescribed angle and measuring the out-of-plane amplitude of the buckled ligament.
The framing is therefore: \emph{calibrate in the accessible in-plane modes,
predict in rotation}, the same constitutive law serving both.

\paragraph{Domain truncation: the grip distance is a physical parameter.}
The RVE truncates the sheet at a window radius $r_{\mathrm{win}}$ around the
ligament, which is legitimate only if the disturbance decays over that distance.
It does not. A convergence study on a $w_{\mathrm{lig}}=18\,\mathrm{mm}$ hinge in
shear gives a peak force that is logarithmic in the truncation radius,
\begin{equation}
  F_s^{\max}\ [\mathrm{N}] \;=\; 303 - 41.2\,\ln r_{\mathrm{win}}\ [\mathrm{mm}],
  \qquad R^2=0.997,
\end{equation}
the signature of a developable cone anchored at the cut tip, whose energy scales
as $B\psi^2\ln(R_{\mathrm{out}}/R_{\mathrm{core}})$ and whose field decays as
$1/r$ --- long-range, with no Saint-Venant plateau. Consequently there is no
``converged'' domain size, and a hinge coupon gripped close to the ligament measures
a stiffer hinge than the same ligament embedded in a large sheet: truncating at
$43\,\mathrm{mm}$ overestimates the force by roughly $60\%$ relative to the
gravity-bounded outer scale. Comparison between experiment and simulation is
therefore only meaningful when the physical grip separation equals the simulated
$r_{\mathrm{win}}$; both are set to $100\,\mathrm{mm}$. In the deployed structure
the outer scale is set instead by the elasto-gravitational length
$\ell_{\mathrm{eg}}=(B/\rho_s g)^{1/3}\approx170\,\mathrm{mm}$, to which the
logarithmic law extrapolates. The same construction underlies the localized-buckling
analysis of Section~\ref{secD}.

\paragraph{Hinge-level validation.}
Single-hinge coupons ($w_{\mathrm{lig}}=18\,\mathrm{mm}$, $\alpha=90^\circ$,
fillet ratio $0.16$, $t=0.5\,\mathrm{mm}$) cut from the same sheet were tested in
fold and in opening. In fold, the measured out-of-plane buckle height tracks the
simulated $u_z(\vartheta)$ through $50^\circ$ with no material tuning --- the
imperfection amplitude being the single adjusted parameter --- and both saturate
near $25\,\mathrm{mm}$. The measurement collapses onto a one-parameter law
\begin{equation}
  u_z(\vartheta) = A\,\sqrt{\sin\vartheta},
  \qquad A\simeq25.7\,\mathrm{mm},\quad R^2=0.97,
  \label{eq:buckle-law}
\end{equation}
which reproduces the $\sqrt{\ }$ post-buckling scaling at small angles and
saturates at large angles as the finite ligament exhausts its developable-cone
reach; a sweep over a fourfold range of fillet ratios leaves the amplitude
$A=29\pm1\,\mathrm{mm}$, so the fillet controls the fracture point but not the
buckling amplitude. In opening the predicted and measured curves share the same
signature --- elastic rise, yield plateau, tearing drop --- with the model
over-predicting the peak force; that comparison is confounded because the same
specimen had previously been folded to $90^\circ$ and was therefore pre-creased,
and is being repeated on virgin coupons. Two caveats bound the calibration as it
stands. The effective ligament section is narrower than nominal, the fillet disc
being centred on the cut tip so that the throat width is
$w_{\mathrm{lig}}(1-\rho_f)$; and the coupons are quasi-static
($1\,\mathrm{mm\,min^{-1}}$) whereas deployment proceeds at centimetres per
second, a regime in which PET stiffens and embrittles, so the slowly-measured
fracture strain is optimistic and is carried with a rate margin.

\section{Results}\label{sec2}
\subsection{Flat-sheet Kirigami}

Material

Calibration of the model

Analysis of motifs of tiles

Actual kirigami cutting


\subsection{Computational efficiency}


\section{Discussion}\label{secD}

\subsection{Locality of out-of-plane buckling of the hinges : comparing paper and steel}

In the thin-hinge regime $\delta\gg t$, the tension-induced instability need not
tilt the square domains out of plane: when they are held coplanar by gravity or an
overburden of areal weight $w=\rho g t+q$, buckling can instead localize in the
ligaments. The selection between this localized mode and the domain-tilting mode,
together with the fold geometry, follows in closed form.

The relative motion of two adjacent domains shares only the primary-cut tip $O$, and
by Chasles' theorem a rigid motion with a fixed point is a rotation about it. That $O$
is fixed follows energetically: translating it strains the junction (cost $\sim Yu^2$,
with $Y=Et$), whereas holding it fixed and folding costs only bending
($\sim B\theta^2/\delta$, with $B=Et^3/[12(1-\nu^2)]$), the former larger by
$(\delta/t)^2\gg1$; symmetry about the fold plane suppresses the remaining out-of-fold
rotations. The fold apex localizes at $O$ rather than in the ligament interior because
a developable fold concentrates an angle deficit $\sim 2\theta$ that, in the interior,
would form a disclination requiring stretching ($\sim Y\theta^2\delta^2$), whereas at
the free edge it is absorbed by the cut opening at no stretching cost---that is, by the
cut-opening angle $2\theta_o$.

The bulge is therefore a developable cone anchored at $O$; with $\delta$ the only
available transverse length, its width is $\lambda\simeq\delta$. The idealized cone's
curvature discontinuities are smoothed over a bending boundary layer
$\ell_B=[B/(Y\theta^2)]^{1/2}$; since the buckling angle
$\theta\simeq\sqrt{2\varepsilon_c}=t/\delta$ (with $\varepsilon_c\simeq\tfrac12(t/\delta)^2$),
one has $\ell_B\sim\delta$ and the fold is a single smooth mound of finite curvature.

Mode selection follows from a moment balance about the hinge axis of one domain.
Gravity resists tilting with $\mathcal{M}_{\mathrm{grav}}\simeq\tfrac12 wL^3$, while the
folded hinge drives it with $\mathcal{M}_{\mathrm{el}}=E I_o\,\theta/\lambda\simeq
Et^4/(12\delta)$ (with $I_o=\delta t^3/12$). Introducing the elasto-gravity length
$\ell_{\mathrm{eg}}=(B/w)^{1/3}$, their ratio is
\begin{equation}
\Gamma\equiv\frac{\mathcal{M}_{\mathrm{grav}}}{\mathcal{M}_{\mathrm{el}}}
=\frac{6\,w\,L^{3}\,\delta}{E\,t^{4}}
=C\left(\frac{L}{\ell_{\mathrm{eg}}}\right)^{3}\frac{\delta}{t},
\qquad C=O(1).
\end{equation}
The domains remain coplanar (localized buckling) for $\Gamma>1$, i.e.\ above a critical
edge length $L^{\star}\sim\ell_{\mathrm{eg}}(t/\delta)^{1/3}$, and tilt below it. The single
inequality $\delta\gg t$ underlies both the pinning of $O$ and the isometry of the fold;
as $\delta\to t$ the junction stretches and peels, $O$ unpins, and the planar auxetic
response of the thick limit is recovered.



\begin{table}[h!]
\centering
\begin{tabular}{lcc}
\toprule
\textbf{Material} & \textbf{Sheet Thickness $t$ (mm)} & \textbf{Critical Length $L^\star$ (cm)} \\
\midrule
Paper & 0.2 & 7.3 \\
Polypropylene & 0.8 & 22.2 \\
Aluminium & 1.0 & 69.4 \\
Soft Steel & 1.5 & 103.7 \\
\bottomrule
\end{tabular}
\caption{Critical lengths $L^\star$ for several thin materials
($\bar{\varepsilon}=1.2$, $\delta=5$~mm).}
\label{tab:longueur_critique}
\end{table}


\section{Conclusion}\label{sec13}

Conclusions may be used to restate your hypothesis or research question, restate your major findings, explain the relevance and the added value of your work, highlight any limitations of your study, describe future directions for research and recommendations. 

In some disciplines use of Discussion or 'Conclusion' is interchangeable. It is not mandatory to use both. Please refer to Journal-level guidance for any specific requirements. 

\backmatter

\bmhead{Supplementary information}

If your article has accompanying supplementary file/s please state so here. 

Authors reporting data from electrophoretic gels and blots should supply the full unprocessed scans for key as part of their Supplementary information. This may be requested by the editorial team/s if it is missing.

Please refer to Journal-level guidance for any specific requirements.


\noindent
If any of the sections are not relevant to your manuscript, please include the heading and write `Not applicable' for that section. 

%%===================================================%%
%% For presentation purpose, we have included        %%
%% \bigskip command. Please ignore this.             %%
%%===================================================%%
\bigskip
\begin{flushleft}%
Editorial Policies for:

\bigskip\noindent
Springer journals and proceedings: \url{https://www.springer.com/gp/editorial-policies}

\bigskip\noindent
Nature Portfolio journals: \url{https://www.nature.com/nature-research/editorial-policies}

\bigskip\noindent
\textit{Scientific Reports}: \url{https://www.nature.com/srep/journal-policies/editorial-policies}

\bigskip\noindent
BMC journals: \url{https://www.biomedcentral.com/getpublished/editorial-policies}
\end{flushleft}

\begin{appendices}

\section{Section title of first appendix}\label{secA1}

An appendix contains supplementary information that is not an essential part of the text itself but which may be helpful in providing a more comprehensive understanding of the research problem or it is information that is too cumbersome to be included in the body of the paper.

%%=============================================%%
%% For submissions to Nature Portfolio Journals %%
%% please use the heading ``Extended Data''.   %%
%%=============================================%%

%%=============================================================%%
%% Sample for another appendix section			       %%
%%=============================================================%%

%% \section{Example of another appendix section}\label{secA2}%
%% Appendices may be used for helpful, supporting or essential material that would otherwise 
%% clutter, break up or be distracting to the text. Appendices can consist of sections, figures, 
%% tables and equations etc.

\end{appendices}

%%===========================================================================================%%
%% If you are submitting to one of the Nature Portfolio journals, using the eJP submission   %%
%% system, please include the references within the manuscript file itself. You may do this  %%
%% by copying the reference list from your .bbl file, paste it into the main manuscript .tex %%
%% file, and delete the associated \verb+\bibliography+ commands.                            %%
%%===========================================================================================%%

\bibliography{sn-bibliography}% common bib file
%% if required, the content of .bbl file can be included here once bbl is generated
%%\input sn-article.bbl

\end{document}
